\subsection{积分的延拓}

由空间 \(L_{F}^{p}\) 的定义可得：具有紧支集的连续函数所成的子空间 \(H_{F}\) 在 \(L_{F}^{p}\) 中稠密（§3, No.~4, 定义 2）。因此，每个连续的（对于 p 阶平均收敛拓扑）线性函数——定义在 \(H_{F}\) 上、取值于一个完备 Hausdorff 拓扑向量空间 G——都可以以唯一方式连续延拓为定义在 \(L_{F}^{p}\) 上、取值于 G 中的连续线性函数（GT, II, §3, No.~6, 定理 2 及 III, §3, No.~1, 命题 3）。

现在，对每个连续函数 \(\mathbf{f}\) （具有紧支集、取值于 Banach 空间 \(\mathrm{F}\) 的），我们已定义过（在 Ch. III, §3, No.~1）它的积分 \(\mu(\mathbf{f}) = \int \mathbf{f} d\mu\) （关于 \(\mu\) 的），它是 \(\mathrm{F}\) 的一个元素；而且我们已证明（Ch. III, §3, No.~2, 命题 6）不等式

\[
\left| \int \mathbf {f} d \mu \right| \leqslant \int | \mathbf {f} | d | \mu | = \mathrm{N} _ {1} (\mathbf {f}).\tag{1}
\]

这个不等式表明线性映射 \(\mathbf{f} \mapsto \int \mathbf{f} d\mu\) （\(\mathcal{H}_{\mathrm{F}}\) 到 \(\mathrm{F}\) 中的）对于 \(\mathcal{H}_{\mathrm{F}}\) 中的平均收敛拓扑是连续的。因此它可以连续延拓到整个空间 \(L_{F}^{1}\) ，从而我们可以给出下述定义：

定义 1. --- 属于 \(\mathcal{L}_{\mathrm{F}}^{1}(\mathrm{X},\mu)\) 的函数称为关于测度 \(\mu\) 可积的（或者再称它们是 \(\mu\) -可积的）。（关于 \(\mu\) 的）可积函数 \(\mathbf{f}\) 的积分，按定义是 \(\mathbf{f}\) 处所取的值，这个值是把一个线性映射连续延拓到 \(\mathcal{L}_{\mathrm{F}}^{1}\) 而得的，所说线性映射即 \(\mathbf{g} \mapsto \int \mathbf{g} d\mu\) （\(\mathcal{K}_{\mathrm{F}}\) 到 \(\mathrm{F}\) 中的）；它仍记作 \(\mu(\mathbf{f})\) 或 \(\int \mathbf{f} d\mu\) ，或 \(\int \mathbf{f}(x) d\mu(x)\) 或 \(\int \mathbf{f}\mu\) ，或 \(\int \mathbf{f}(x)\mu(x)\) 。

例子. --- 设 X 是一个离散空间，\(\mu\) 是 X 上的一个测度，并令 \(\alpha(x) = \mu(\varphi_{\{x\}})\) 对每个 \(x \in X\) 。此时 \(F_{F}^{1}\) 中的函数都是可积的，换句话说 \(L_{F}^{1} = F_{F}^{1}\) ；此外，对每个函数 \(f \in L_{F}^{1}\) ，

\[
\int \mathbf {f} d \mu = \sum_ {x \in X} \alpha (x) \mathbf {f} (x).
\]

因为，设 \(\mathbf{f} \in \mathcal{F}_F^1\) ；我们有 \(|\mu|^*(|\mathbf{f}|) = \sum_{x \in X} |\alpha(x)| \cdot |\mathbf{f}(x)| < +\infty\) （§1, No.~3, 例子）；于是对每个 \(\varepsilon > 0\) ，存在一个有限子集 \(M\) （\(X\) 的）使得

\[
\sum_ {x \in \mathrm{X} - \mathrm{M}} | \alpha (x) | \cdot | \mathbf {f} (x) | \leqslant \varepsilon .
\]

函数 \(\mathbf{g}\) 与 \(\mathbf{f}\) 在诸点 \(x\in M\) （\(|\mathbf{f}|\) 在那里为有限）处相等，在别处等于 0，它属于 \(\mathcal{H}(\mathrm{X};\mathrm{F})\) ，而且按已作的约定，

\[
| \mu | ^ {*} (| \mathbf {f} - \mathbf {g} |) \leqslant \sum_ {\boldsymbol {x} \in \mathrm{X} - \mathrm{M}} | \alpha (\boldsymbol {x}) | \cdot | \mathbf {f} (\boldsymbol {x}) | \leqslant \varepsilon ,
\]

这就证明了 \(\mathbf{f} \in \mathcal{L}_{\mathrm{F}}^{1}\) 。另一方面，

\[
\left| \mu (\mathbf {g}) - \sum_ {x \in X} \alpha (x) \mathbf {f} (x) \right| \leqslant \sum_ {x \in X - M} | \alpha (x) | \cdot | \mathbf {f} (x) | \leqslant \varepsilon ,
\]

由此得第二个断言。

换句话说，\(\mu\) -可积函数 \(\mathbf{f}\) 就是使族 \((\alpha(x)\mathbf{f}(x))_{x\in X}\) 绝对可和（GT, IX, §3, No.~6）的那些函数，且积分 \(\int \mathbf{f} d\mu\) 是这个族的和。

由于 \(\mu(\mathbf{f})\) 依定义在 \(\mathcal{L}_{\mathrm{F}}^{1}\) 上连续，且它取值于一个 Hausdorff 空间，故 \(\mu(\mathbf{f}) = 0\) 对每个属于 0 在 \(\mathcal{L}_{\mathrm{F}}^{1}\) 中的闭包的函数（即可忽略的函数）成立；若 \(\mathbf{f}\) 与 \(\mathbf{g}\) 是两个等价的可积函数，则 \(\mu(\mathbf{f}) = \mu(\mathbf{g})\) 。换句话说，\(\mu(\mathbf{f})\) 的值只依赖于类 \(\widetilde{\mathbf{f}}\) （可积函数 \(\mathbf{f}\) 的）；它仍记作 \(\mu(\widetilde{\mathbf{f}})\) ，而函数 \(\widetilde{\mathbf{f}} \mapsto \mu(\widetilde{\mathbf{f}})\) 是 \(L_{F}^{1}\) 到 F 中的一个连续线性映射。若一个取值于 F、在 X 中几乎处处有定义的函数 f 等价于一个可积函数，我们仍说 f 是可积的，并写 \(\int f d\mu = \mu(\widetilde{\mathbf{f}})\) ；类似地可定义取值于 \(\overline{R}\) 、几乎处处有定义且有限的可积函数及其积分。

